\chapter{Vasectomy}

\section*{Overview}
A vasectomy is a surgical procedure for male sterilization or permanent contraception. The vas deferens are severed and sealed to prevent sperm from entering the seminal fluid. The exact approach, setting, and preparation depend on the indication, anatomy, and the patient's health.

\section*{Alternative Names}
Male sterilization

\section*{Principle}
Each vas deferens is identified, isolated through a small scrotal incision, and a segment is cut and sealed with ligation or cautery. Blocking the sperm transport pathway means the ejaculate contains seminal fluid without sperm.
\section*{History}
First performed on dogs by Sir Astley Cooper in 1830. Regarded as a sterilization method in the early 20th century. The 'no-scalpel' technique was developed in China in 1974.

\section*{Why It Is Done}
Ordered by men seeking permanent, highly effective contraception.

\section*{Is This a Primary or Secondary Test or Procedure}
Primary, highly effective method of permanent male contraception.

\section*{How It Is Performed}
Usually performed under local anesthesia in an office setting. The clinician makes small punctures in the scrotum, isolates the vas deferens, cuts them, and seals the ends with clips or cautery.

\section*{Preparation}
Shave the scrotum and clean the area on the day of the procedure. Wear tight-fitting underwear or a jockstrap to support the scrotum afterward.

\section*{Contraindications and Precautions}
Active scrotal infection, local anatomical abnormalities (large varicocele, inguinal hernia), or uncertainty about permanent sterilization.

\section*{Risks}
Hematoma, infection, sperm granuloma, chronic post-vasectomy pain syndrome (<1%), or rare failure/recanalization (<0.1%).

\section*{Aftercare and Recovery}
Apply ice packs to the scrotum. Rest for 24-48 hours. Avoid heavy lifting and sexual activity for 1 week. Use back-up contraception until cleared.

\section*{Outcome and Follow-up}
A semen analysis is performed 8-12 weeks post-procedure. Sterilization is confirmed only when the semen shows azoospermia.

\section*{Expected Results}
Successful occlusion of the vas deferens; post-vasectomy semen analysis demonstrating zero sperm (azoospermia).

\section*{Follow-up}
Semen analysis at 12 weeks or after 20 ejaculations.

\section*{When to Seek Medical Care}
Follow the procedure team's specific instructions. Seek urgent medical assessment for severe or worsening pain, uncontrolled bleeding, fever or signs of infection, trouble breathing, chest pain, fainting, new weakness or confusion, inability to urinate, or any rapidly worsening symptom. The appropriate warning signs depend on the procedure and the patient's medical condition.

\section*{References}
\begin{enumerate}
  \item MedlinePlus. Vasectomy. U.S. National Library of Medicine; 2025.
  \item Current Medical Diagnosis and Treatment. McGraw-Hill; 2025.
\end{enumerate}

\clearpage
